% GENERATED FILE. Edit canonical lessons, then run npm run generate:books.
% canonical-source-hash: fnv1a64:8cb27bba6e73d17a
% canonical-lessons: AR-C85-khafif, AR-C85-thaqil, AR-C85-sahih, AR-C85-khata, AR-C85-muhimm

\chapter{Light, Heavy, Correct, Mistake, Important}
\label{ch:ar-light-heavy-correct-mistake-important}
\hlchaptermodality{85}

\begin{chapteropening}
\textbf{By the end of this chapter:} \emph{I can say light, heavy, correct, true, a mistake and important.}

Complete the last lesson of chapter 85: I can say light, heavy, correct, true, a mistake and important.
\end{chapteropening}

\section[\texorpdfstring{\ar{خفيف} (khafīf)}{خفيف (khafīf)}]{\ar{خفيف} (khafīf) — light (in weight)}

\label{lesson:AR-C85-khafif}



\begin{quote}
Before the new one: say the Arabic for clean, then the Arabic for dirty.
\end{quote}

\subsection*{You'll want to know: \ar{خفيف}}
\textbf{\ar{خفيف}} — \emph{khafīf} — \textquotedblleft{}light (in weight)\textquotedblright{}.

\begin{cousinweb}
\textbf{\ar{خفيف}} (\emph{khafīf}) is built on the native Arabic root \emph{kh-f-f}, \textquotedblleft{}to be light\textquotedblright{} (Lane). \emph{Khafīf} is the \emph{faʿīl} adjective, \textquotedblleft{}light\textquotedblright{}.
\end{cousinweb}

\begin{grammarlens}[title={What you've built}]
One more word for this chapter.
\end{grammarlens}

\subsection*{Your turn}
\begin{itemize}
\raggedright
  \item \emph{Say it:} \emph{khafīf}
  \item \emph{Say it:} \emph{khafīf}, once more
  \item \emph{Say it:} say \emph{khafīf}
  \item \emph{Recall:} say the Arabic for clean, then the Arabic for dirty, then say \emph{khafīf} again
\end{itemize}

\begin{tcolorbox}[breakable,colback=teal!4,colframe=teal!35!black,title={Before you move on}]
What does \ar{خفيف} mean? (\textquotedblleft{}Light (in weight)\textquotedblright{}.) Say it once more.
\end{tcolorbox}
\section[\texorpdfstring{\ar{ثقيل} (thaqīl)}{ثقيل (thaqīl)}]{\ar{ثقيل} (thaqīl) — heavy}

\label{lesson:AR-C85-thaqil}



\begin{quote}
Before the new one: say the Arabic for dirty, then the Arabic for light.
\end{quote}

\subsection*{You'll want to know: \ar{ثقيل}}
\textbf{\ar{ثقيل}} — \emph{thaqīl} — \textquotedblleft{}heavy\textquotedblright{}.

\begin{cousinweb}
\textbf{\ar{ثقيل}} (\emph{thaqīl}) is built on the native Arabic root \emph{th-q-l}, \textquotedblleft{}to be heavy\textquotedblright{} (Lane). Hebrew \emph{shaqal}, \textquotedblleft{}to weigh\textquotedblright{}, the source of English \emph{shekel}, is usually counted a cousin.
\end{cousinweb}

\begin{grammarlens}[title={What you've built}]
One more word for this chapter.
\end{grammarlens}

\subsection*{Your turn}
\begin{itemize}
\raggedright
  \item \emph{Say it:} \emph{thaqīl}
  \item \emph{Say it:} \emph{thaqīl}, once more
  \item \emph{Say it:} say \emph{thaqīl}
  \item \emph{Recall:} say the Arabic for dirty, then the Arabic for light, then say \emph{thaqīl} again
\end{itemize}

\begin{tcolorbox}[breakable,colback=teal!4,colframe=teal!35!black,title={Before you move on}]
What does \ar{ثقيل} mean? (\textquotedblleft{}Heavy\textquotedblright{}.) Say it once more.
\end{tcolorbox}
\section[\texorpdfstring{\ar{صحيح} (ṣaḥīḥ)}{صحيح (ṣaḥīḥ)}]{\ar{صحيح} (ṣaḥīḥ) — correct, true}

\label{lesson:AR-C85-sahih}



\begin{quote}
Before the new one: say the Arabic for light, then the Arabic for heavy.
\end{quote}

\subsection*{You'll want to know: \ar{صحيح}}
\textbf{\ar{صحيح}} — \emph{ṣaḥīḥ} — \textquotedblleft{}correct, true\textquotedblright{}.

\begin{cousinweb}
\textbf{\ar{صحيح}} (\emph{ṣaḥīḥ}) is built on the native Arabic root \emph{ṣ-ḥ-ḥ}, \textquotedblleft{}to be sound, free from defect\textquotedblright{} (Lane). \emph{Ṣaḥīḥ} is the \emph{faʿīl} adjective, \textquotedblleft{}sound, correct\textquotedblright{}.
\end{cousinweb}

\begin{grammarlens}[title={What you've built}]
One more word for this chapter.
\end{grammarlens}

\subsection*{Your turn}
\begin{itemize}
\raggedright
  \item \emph{Say it:} \emph{ṣaḥīḥ}
  \item \emph{Say it:} \emph{ṣaḥīḥ}, once more
  \item \emph{Say it:} say \emph{ṣaḥīḥ}
  \item \emph{Recall:} say the Arabic for light, then the Arabic for heavy, then say \emph{ṣaḥīḥ} again
\end{itemize}

\begin{tcolorbox}[breakable,colback=teal!4,colframe=teal!35!black,title={Before you move on}]
What does \ar{صحيح} mean? (\textquotedblleft{}Correct, true\textquotedblright{}.) Say it once more.
\end{tcolorbox}
\section[\texorpdfstring{\ar{خطأ} (khaṭaʾ)}{خطأ (khaṭaʾ)}]{\ar{خطأ} (khaṭaʾ) — a mistake; wrong}

\label{lesson:AR-C85-khata}



\begin{quote}
Before the new one: say the Arabic for heavy, then the Arabic for correct.
\end{quote}

\subsection*{You'll want to know: \ar{خطأ}}
\textbf{\ar{خطأ}} — \emph{khaṭaʾ} — \textquotedblleft{}a mistake; wrong\textquotedblright{}.

\begin{cousinweb}
\textbf{\ar{خطأ}} (\emph{khaṭaʾ}) is built on the native Arabic root \emph{kh-ṭ-ʾ}, \textquotedblleft{}to err, to sin\textquotedblright{} (Lane). Hebrew \emph{ḥeṭ}, \textquotedblleft{}sin\textquotedblright{}, is usually counted a cousin.
\end{cousinweb}

\begin{grammarlens}[title={What you've built}]
One more word for this chapter.
\end{grammarlens}

\subsection*{Your turn}
\begin{itemize}
\raggedright
  \item \emph{Say it:} \emph{khaṭaʾ}
  \item \emph{Say it:} \emph{khaṭaʾ}, once more
  \item \emph{Say it:} say \emph{khaṭaʾ}
  \item \emph{Recall:} say the Arabic for heavy, then the Arabic for correct, then say \emph{khaṭaʾ} again
\end{itemize}

\begin{tcolorbox}[breakable,colback=teal!4,colframe=teal!35!black,title={Before you move on}]
What does \ar{خطأ} mean? (\textquotedblleft{}A mistake; wrong\textquotedblright{}.) Say it once more.
\end{tcolorbox}
\section[\texorpdfstring{\ar{مهم} (muhimm)}{مهم (muhimm)}]{\ar{مهم} (muhimm) — important}

\label{lesson:AR-C85-muhimm}



\begin{quote}
Before the new one: say the Arabic for correct, then the Arabic for a mistake; wrong.
\end{quote}

\subsection*{You'll want to know: \ar{مهم}}
\textbf{\ar{مهم}} — \emph{muhimm} — \textquotedblleft{}important\textquotedblright{}.

\begin{cousinweb}
\textbf{\ar{مهم}} (\emph{muhimm}) is built on the native Arabic root \emph{h-m-m}, \textquotedblleft{}to concern, to make anxious\textquotedblright{} (Lane). \emph{Muhimm} is the Form IV participle, \textquotedblleft{}that which concerns\textquotedblright{}.
\end{cousinweb}

\begin{grammarlens}[title={What you've built}]
One more word for this chapter.
\end{grammarlens}

\subsection*{Your turn}
\begin{itemize}
\raggedright
  \item \emph{Say it:} \emph{muhimm}
  \item \emph{Say it:} \emph{muhimm}, once more
  \item \emph{Say it:} say \emph{muhimm}
  \item \emph{Recall:} say the Arabic for correct, then the Arabic for a mistake; wrong, then say \emph{muhimm} again
\end{itemize}

\begin{tcolorbox}[breakable,colback=teal!4,colframe=teal!35!black,title={Before you move on}]
What does \ar{مهم} mean? (\textquotedblleft{}Important\textquotedblright{}.) Say it once more.
\end{tcolorbox}
